\section{习题与解答}
\begin{exercise}\label{ex:m1anisotropicsource}
光源静止系单位固有时、单位立体角辐射能为$I'(\theta')=I_0(1+a\cos\theta')$，$|a|\leq1$。源沿对称轴以$V$运动。求实验室发射时间角功率和总功率，讨论何时总功率大于静止系值。
\Diagram{fig:m1anisotropicq}{源系辐射含沿运动轴的一阶不对称项}{\draw[axis](-2,0)--(3.1,0)node[right]{$x',V$};\draw[curve,domain=0:360,samples=100]plot({1.2*(1+.6*cos(\x))*cos(\x)},{1.2*(1+.6*cos(\x))*sin(\x)});\node[above]at(0,1.8){图取$a>0$};}
\end{exercise}
\begin{solution}
图\ref{fig:m1anisotropicq}总静止功率仍为$P_0=4\pi I_0$，因为$\cos\theta'$全空间积分为零。对同一批光子使用能量、时间与立体角三因子，
\[
 I(\theta)=\frac{\delta^3}{\gamma}I_0
 \left[1+a\frac{\cos\theta-\beta}{1-\beta\cos\theta}\right].
\]
要积分总功率，改在源系按$\mu'=\cos\theta'$积分更直接。每束的实验室能量乘$\gamma(1+\beta\mu')$，除发射时间$\gamma\dd\tau$，所以
\begin{align*}
 P&=2\pi I_0\int_{-1}^1(1+a\mu')(1+\beta\mu')\dd\mu'\\
 &=2\pi I_0\int_{-1}^1[1+(a+\beta)\mu'+a\beta\mu'^2]\dd\mu'\\
 &=2\pi I_0\left[\mu'+\frac{a+\beta}{2}\mu'^2+
 \frac{a\beta}{3}\mu'^3\right]_{-1}^{1}
 =4\pi I_0(1+a\beta/3).
\end{align*}
\Diagram{fig:m1anisotropica}{总功率的额外项来自源系净辐射动量}{\draw[curve](0,0)rectangle(2.7,1);\node at(1.35,.5){$P_0$};\draw[->](2.9,.5)--(3.7,.5);\node[right]at(3.9,.5){$P=P_0+\beta c\dot p'_x$};\node[below]at(2,-.7){$\dot p'_x=aP_0/(3c)$};}
图\ref{fig:m1anisotropica}中，总功率的修正项由源系净辐射动量决定。$a\beta>0$时实验室功率更大，反号时更小，$a=0$才恢复各向同性源的总功率不变。给定发射机制若长期维持不对称辐射，光源会有反冲；本题假定外部装置维持其匀速运动，或只研究一个瞬时发射小过程。
\end{solution}

\begin{exercise}[选做]\label{ex:m1antitelephone}
假设存在在任意发射者静止系均可沿任意方向以同一速率$u>c$发送的信号。地球$A$固定在$x=0$，飞船$B$满足$x=L+Vt$。$A$在$t=0$发信号，$B$收到后立即在自身系以$-u$回复。求回复到达地球时刻，讨论因果矛盾的条件。
\Diagram{fig:m1antiq}{假设的超光速信号是类空路径；此图不表示真实光线}{\draw[axis](-.3,0)--(3.7,0)node[right]{$x/L$};\draw[axis](0,-1.2)--(0,3.4)node[above]{$ct/L$};\draw[curve](1,0)--(3.7,3);\draw[->,very thick](0,0)--(1.818182,.909091);\draw[->,very thick](1.818182,.909091)--(0,-.413223);\node[right]at(1.818182,.909091){$R$};\node[left]at(0,-.413223){$C$};}
\end{exercise}
\begin{solution}
出发信号$x=ut$与飞船相交，图\ref{fig:m1antiq}的接收事件为$t_R=L/(u-V)$、$x_R=uL/(u-V)$。以飞船系中从回复发出起经过的坐标时间$s>0$作参数，回复有$\Delta t'=s$、$\Delta x'=-us$。逆变换逐项给
\[
 \Delta t=\gamma(1-Vu/c^2)s,\qquad
 \Delta x=\gamma(V-u)s.
\]
回复到达地球要求$x_R+\Delta x=0$，所以$s=x_R/[\gamma(u-V)]>0$。代入时间增量，
\[
 t_C=t_R+\frac{(1-Vu/c^2)x_R}{u-V}
 =\frac{L[2u-V(1+u^2/c^2)]}{(u-V)^2}.
\]
当$Vu=c^2$，回复路径在实验室的时间增量为零，仍由以上参数式处理；直接以$\Delta x/\Delta t$表示斜率则会出现零分母。当$Vu>c^2$，沿回复路径的实验室时间增量为负，不能再用面向未来的普通速度图像解释其方向。
\Diagram{fig:m1antia}{若回复早于出发，地球可在发送之前收到由该发送触发的回信}{\draw[axis](0,-1.4)--(0,2)node[above]{地球钟};\fill(0,0)circle(1.5pt)node[right]{$A:t=0$发送};\fill(0,-.9)circle(1.5pt)node[right]{$C:t_C<0$收到回复};\draw[->,thick](-.5,-.8)--(-.5,-.1);}
因此$V>2u/(1+u^2/c^2)$时$t_C<0$。对$u>c$，这个门槛小于$c$，可由亚光速飞船达到。图\ref{fig:m1antiq}取$u=2c,V=.9c$作示意，回复位于地球出发事件的过去。图\ref{fig:m1antia}所示矛盾依赖“每个惯性系都能双向按同一超光速规则发信号”的假设；视超光速图样不具有这种信号能力，不能用它构造本装置。
\end{solution}

\begin{exercise}\label{ex:m1clocktrain}
列车固有长$L$，以速度$V$运动。车头和车尾各有一只在车系同步的钟。地面在$t=0$同时读取两钟，求车头减车尾的读数差，并求地面收到这两次读数的先后是否必然相同。
\Diagram{fig:m1trainq}{同时读取运动列车的两端钟}{\draw[axis](-.5,0)--(5,0)node[right]{$x$};\draw[curve](0,.5)rectangle(3.5,1.3);\node[above]at(0,1.3){尾};\node[above]at(3.5,1.3){头};\draw[axis](3.8,.9)--(4.8,.9)node[above]{$V$};\draw[guide](0,2)--(3.5,2);\node[above]at(1.75,2){$t=0$};}
\end{exercise}
\begin{solution}
图\ref{fig:m1trainq}中的同时截面长为$L/\gamma$。用$\Delta t'=\gamma(\Delta t-V\Delta x/c^2)$，得到头钟减尾钟
\[
 \Delta t'=-\gamma\frac{V}{c^2}\frac{L}{\gamma}=-\frac{VL}{c^2}.
\]
头钟在地面同时截面上落后。读数事件之间是类空间隔，其光信号到达同一观察者的时间还要加各自传播时间。
\Diagram{fig:m1traina}{地面接收者的位置影响读数信号的到达次序}{\draw[curve](0,0)--(4,0);\fill(0,0)circle(1.4pt)node[below]{尾};\fill(4,0)circle(1.4pt)node[below]{头};\fill(1,1.5)circle(1.4pt)node[above]{$D$};\draw[->](0,0)--(1,1.5);\draw[->](4,0)--(1,1.5);\node[below]at(2,-.5){$\Delta t_{\rm arr}=(r_{\text{头}}-r_{\text{尾}})/c$};}
图\ref{fig:m1traina}中的接收者若更靠近车尾，先收到尾部信号；靠近车头则相反；在垂直平分线上同时收到。钟面读数的差与信号到达次序没有一一对应关系。
\end{solution}

\begin{exercise}\label{ex:m1lifetime}
不稳定粒子的平均固有寿命为$\tau_0$，在实验室以$v$匀速飞行。求在距离$L$处仍未衰变的比例。再在粒子系说明为什么得到同一结果。
\Diagram{fig:m1lifeq}{粒子从产生点飞向距离为$L$的探测平面}{\draw[axis](-.3,0)--(5,0)node[right]{$x$};\fill(0,0)circle(2pt)node[below]{$O$};\draw[curve](4,-.7)--(4,1.2)node[above]{探测面};\draw[->,thick](.4,.3)--(2,.3)node[above]{$v$};\draw[<->](0,-.5)--(4,-.5)node[midway,below]{$L$};}
\end{exercise}
\begin{solution}
图\ref{fig:m1lifeq}中实验室飞行时间为$t=L/v$，粒子经历的固有时为$\tau=t/\gamma_v$。恒定衰变概率对应$\dd N/\dd\tau=-N/\tau_0$。分离变量并用$N(0)=N_0$，
\[
 \int_{N_0}^{N(\tau)}\frac{\dd n}{n}
 =-\frac1{\tau_0}\int_0^\tau\dd s,
 \qquad \ln\frac{N(\tau)}{N_0}=-\frac{\tau}{\tau_0}.
\]
取指数并代入飞行固有时，得
\[
 \frac N{N_0}=\exp\left[-\frac{L}{\gamma_vv\tau_0}\right].
\]
\Diagram{fig:m1lifea}{粒子系中探测面向静止粒子运动}{\fill(0,0)circle(2pt)node[above]{粒子};\draw[curve](3,-.7)--(3,1.2);\draw[axis](2.7,.5)--(1.5,.5)node[midway,above]{$v$};\draw[<->](0,-.5)--(3,-.5)node[midway,below]{$L/\gamma_v$};}
在图\ref{fig:m1lifea}的粒子系，产生点到探测面在同一时刻的距离缩短为$L/\gamma_v$，探测面以$v$靠近，所需时间$L/(\gamma_vv)$，恰为上面的固有时间。这里“平均寿命”不代表所有粒子都在同一时刻衰变，结果是统计存活比例。
\end{solution}

\begin{exercise}\label{ex:m1crossvelocity}
飞船$S'$以$V$相对实验室运动，在船上沿$y'$方向发射速度为$u$的粒子。求实验室速度、洛伦兹因子及粒子运动到船系坐标$y'=H$时的实验室位置。
\Diagram{fig:m1crossq}{船内的横向发射}{\draw[curve](0,0)rectangle(3,2.2);\fill(1,0)circle(1.5pt);\draw[axis](1,0)--(1,1.8)node[left]{$u$};\node[right]at(1,2){$H$};\draw[axis](3.4,1)--(4.6,1)node[right]{$V$};}
\end{exercise}
\begin{solution}
图\ref{fig:m1crossq}中$v'_x=0,v'_y=u$。逆速度变换给$v_x=V,v_y=u/\gamma$。于是
\[
 1-v^2/c^2=(1-\beta^2)(1-u^2/c^2),\quad \gamma_v=\gamma\gamma_u.
\]
到达上边界的事件在船系为$t'=H/u,x'=0,y'=H$。逆坐标变换得$t=\gamma H/u$，$x=\gamma VH/u$，$y=H$。
\Diagram{fig:m1crossa}{实验室中粒子轨迹斜向前方，与船保持相同纵向速度}{\draw[axis](-.3,0)--(4.8,0)node[right]{$x$};\draw[axis](0,-.2)--(0,2.5)node[above]{$y$};\draw[very thick,->](0,0)--(3.7,2);\draw[guide](0,2)--(3.7,2)--(3.7,0);\node[left]at(0,2){$H$};\node[below]at(3.7,0){$\gamma VH/u$};}
图\ref{fig:m1crossa}的斜率$H/(\gamma VH/u)=u/(\gamma V)$也等于$v_y/v_x$。若直接使用伽利略合成$v_y=u$，会使坐标变换与轨迹速度互相矛盾。
\end{solution}

\begin{exercise}\label{ex:m1coneenergy}
静止系各向同性的单次闪光总能量为$E_0$，各方向的光谱相同。实验室观察前向锥$\theta\leq\arccos\beta$。求锥内光子数与能量各占整次闪光的比例。
\Diagram{fig:m1coneq}{前向锥恰好由静止系前半球变换而来}{\draw[axis](-.3,0)--(4,0)node[right]{$x$};\draw[curve](0,0)--(3,1.3);\draw[curve](0,0)--(3,-1.3);\draw[guide](3,-1.3)--(3,1.3);\node[right]at(3,1.3){$\cos\theta_c=\beta$};}
\end{exercise}
\begin{solution}
由例\ref{ex:m1aberration}，图\ref{fig:m1coneq}中的光锥来自$0\leq\mu'\leq1$。静止系光子数各向同性，变换不改变光子数，所以数目占比为$1/2$。
\Diagram{fig:m1conea}{在静止系积分能量权重，权重为$\gamma(1+\beta\mu')$}{\draw[axis](0,0)--(3.8,0)node[right]{$\mu'$};\draw[axis](0,0)--(0,2.8)node[above]{能量权重};\draw[curve](0,1)--(3,2.3);\draw[guide](3,0)--(3,2.3);\node[below]at(3,0){$1$};\node[left]at(0,1){$\gamma$};}
能量按图\ref{fig:m1conea}的权重变换，锥内能量为
\[
 E_c=\frac{\gamma E_0}{2}\int_0^1(1+\beta\mu')\dd\mu'
 =\frac{\gamma E_0}{2}\left[\mu'+\frac\beta2\mu'^2\right]_0^1
 =\gamma E_0\left(\frac12+\frac\beta4\right).
\]
整个闪光的能量为$\gamma E_0$，故能量份额为$1/2+\beta/4$。即使高相对论速度，静止系前半球所含光子的实验室能量份额也只趋向$3/4$；这与小立体角内较大的角能量密度并不冲突。
\end{solution}

\begin{exercise}\label{ex:m1transverse}
运动光源的静止频率为$\nu_0$。分别求实验室观察方向与速度垂直、以及光源系发射方向与速度垂直时的实验室频率。
\Diagram{fig:m1transq}{“横向”可指两个不同惯性系中的方向}{\draw[axis](0,0)--(3.4,0)node[right]{$V$};\draw[axis](0,0)--(0,2)node[above]{$\theta=\pi/2$};\draw[->,thick](0,0)--(2.6,1.6)node[right]{$\theta'=\pi/2$的光};}
\end{exercise}
\begin{solution}
图\ref{fig:m1transq}中竖直的实验室光线有$\cos\theta=0$，所以$\nu=\nu_0/\gamma$，发生红移。若取源系$\cos\theta'=0$，则由逆能量变换$\nu=\gamma\nu_0$，发生蓝移。
\Diagram{fig:m1transa}{两个横向条件对应不同实验室角度}{\draw[axis](0,0)--(3.5,0)node[right]{$x$};\draw[curve](0,0)--(0,2);\draw[curve](0,0)--(2.6,1.6);\node[left]at(0,1.7){$\nu_0/\gamma$};\node[right]at(2.6,1.6){$\gamma\nu_0$};\draw[->](.9,0)arc(0:31.6:.9);\node[above]at(1.3,1.1){$\cos\theta=\beta$};}
第二种光线在实验室满足$\cos\theta=\beta$，如图\ref{fig:m1transa}。将其代入$\nu=\nu_0/[\gamma(1-\beta\cos\theta)]$，分母成为$\gamma(1-\beta^2)=1/\gamma$，与直接变换吻合。叙述横向多普勒效应时必须说明横向角在哪个系定义。
\end{solution}

\begin{exercise}\label{ex:m1fourwork}
静质量$m$的粒子由静止开始受恒定三力$f$沿$x$加速。求其动量、能量、速度和位移随实验室时间的变化，以及经历的固有时。
\Diagram{fig:m1workq}{实验室恒定三力的直线运动}{\draw[axis](-.3,0)--(4.7,0)node[right]{$x$};\fill(0,0)circle(2pt)node[below]{$m$};\draw[very thick,->](.2,.5)--(2,.5)node[midway,above]{$f$};}
\end{exercise}
\begin{solution}
以图\ref{fig:m1workq}的出发事件为$t=x=0$。恒三力满足$\dd p/\dd t=f$；用$p(0)=0$，有$\int_0^{p(t)}\dd p_1=f\int_0^t\dd t_1$，即$p=ft$。质壳关系和$v=pc^2/E$给
\[
 E=mc^2\sqrt{1+s^2},\qquad
 v=\frac{cs}{\sqrt{1+s^2}},\qquad s=\frac{ft}{mc}.
\]
位移的定积分为
\begin{align*}
 x(t)&=\int_0^t v(t_1)\dd t_1
 =\frac{mc^2}f\int_0^s\frac{q\dd q}{\sqrt{1+q^2}}\\
 &=\frac{mc^2}f[\sqrt{1+q^2}]_0^s
 =\frac{mc^2}f[\sqrt{1+s^2}-1].
\end{align*}
于是$fx=E-mc^2$，外力做功恰等于动能增加。
\Diagram{fig:m1worka}{恒力世界线渐近于光线，虚线斜率为1}{\draw[axis](-.2,0)--(4.5,0)node[right]{$fx/(mc^2)$};\draw[axis](0,-.2)--(0,4.5)node[above]{$ft/(mc)$};\draw[curve,domain=0:4,samples=80]plot({sqrt(1+\x*\x)-1},\x);\draw[guide](0,1)--(3.4,4.4);}
固有时间应沿图\ref{fig:m1worka}的加速世界线累计。由$\gamma_v=\sqrt{1+s^2}$，
\[
 \tau=\frac{mc}f\int_0^s\frac{\dd q}{\sqrt{1+q^2}}
 =\frac{mc}f\operatorname{arsinh}s.
\]
最后一步令$q=\sinh\eta$，则$\dd q=\cosh\eta\dd\eta$、$\sqrt{1+q^2}=\cosh\eta$，上下限为$0,\operatorname{arsinh}s$，积分为$[\eta]_0^{\operatorname{arsinh}s}$。反解$s=\sinh(f\tau/mc)$，代回得到
\[
 ct=\frac{mc^2}f\sinh\frac{f\tau}{mc},\qquad
 x=\frac{mc^2}f\left(\cosh\frac{f\tau}{mc}-1\right).
\]
从$v(t)=cs(1+s^2)^{-1/2}$用乘积与链式法则求导，
\[
 \frac{\dd v}{\dd t}=c\frac{\dd s}{\dd t}
 \left[(1+s^2)^{-1/2}-s^2(1+s^2)^{-3/2}\right]
 =\frac fm(1+s^2)^{-3/2}.
\]
纵向固有加速度为$\gamma_v^3\dd v/\dd t=f/m$。恒三力对应恒纵向固有加速度，实验室加速度则不断减小。$s\ll1$时$x\simeq ft^2/(2m)$、$\tau\simeq t$，任何有限$t$都有$v<c$。
\end{solution}
\begin{exercise}\label{ex:m1twoemitters}
静止系中两灯位于$\pm L/2$，同时向中点发光。求两束光相遇事件在运动系的坐标，以及两发射事件在运动系是否同时。
\Diagram{fig:m1emitq}{两灯同时发射相向光线}{\draw[axis](-2.3,0)--(2.5,0);\fill(-2,0)circle(1.5pt)node[below]{$A$};\fill(2,0)circle(1.5pt)node[below]{$B$};\draw[->](-1.8,.5)--(-.2,.5);\draw[->](1.8,.5)--(.2,.5);}
\end{exercise}
\begin{solution}
在图\ref{fig:m1emitq}的$S$系中，左右发射事件分别为$A:(0,-L/2)$、$B:(0,L/2)$，坐标次序为$(ct,x)$。两束光到中点均走$L/2$，故相遇于$C:(L/2,0)$。变换至沿$+x$以$V$运动的$S'$，得到
\begin{align*}
 A'&:(\gamma\beta L/2,-\gamma L/2),\\
 B'&:(-\gamma\beta L/2,\gamma L/2),\\
 C'&:(\gamma L/2,-\gamma\beta L/2).
\end{align*}
所以右灯先发光，$t'_B-t'_A=-\gamma\beta L/c$。
\Diagram{fig:m1emita}{运动系中先发射的光走得更久，相遇点仍唯一}{\draw[axis](-2.7,0)--(2.5,0)node[right]{$x'$};\draw[axis](0,-1.7)--(0,3)node[above]{$ct'$};\draw[very thick](-2,1.2)--(-1.2,2)--(2,-1.2);\node[left]at(-2,1.2){$A$};\node[right]at(2,-1.2){$B$};\node[above]at(-1.2,2){$C$};}
图\ref{fig:m1emita}中，左行和右行光仍具有光速。由$A'$到$C'$，
\[
 c(t'_C-t'_A)=\frac{\gamma L}2(1-\beta)
 =x'_C-x'_A;
\]
由$B'$到$C'$，
\[
 c(t'_C-t'_B)=\frac{\gamma L}2(1+\beta)
 =x'_B-x'_C.
\]
两条光路的传播时间不同，正好补偿了两次发射的不同步。相遇是同一时空事件，不会因参考系改变而消失。$x'_B-x'_A=\gamma L$是两个不同时发射事件的位置差；两灯在$S'$的同时距离仍为$L/\gamma$。
\end{solution}

\begin{exercise}\label{ex:m1clockoffset}
两个静止于$S'$的钟相距$L$，钟面相对坐标时间偏移分别为$b_A,b_B$。实验室在$t=0$同时读取，求读数差；什么偏移能使这次读数相同？
\Diagram{fig:m1offsetq}{钟面读数含可人为设置的常数偏移}{\draw[curve](0,0)circle(.5);\draw[curve](4,0)circle(.5);\node at(0,0){$A$};\node at(4,0){$B$};\node[above]at(0,.7){$t'+b_A$};\node[above]at(4,.7){$t'+b_B$};\draw[axis](1,-1)--(2.5,-1)node[right]{$V$};}
\end{exercise}
\begin{solution}
在图\ref{fig:m1offsetq}的实验室同时截面，两钟坐标时间差$\Delta t'=-VL/c^2$，因此钟面差$T_B-T_A=b_B-b_A-VL/c^2$。使本次读数相同需$b_B-b_A=VL/c^2$。
\Diagram{fig:m1offseta}{两钟速率相同，偏移只改变读数的常数差}{\draw[axis](0,0)--(4.5,0)node[right]{$t$};\draw[axis](0,0)--(0,3.6)node[above]{读数};\draw[curve](0,.4)--(3.8,2.8);\draw[guide](0,1)--(3.8,3.4);\node[right]at(3.8,2.8){固定偏移};}
图\ref{fig:m1offseta}中两钟在实验室的速率均为$1/\gamma$，所以在每个实验室同时截面上，该读数差都相同；这个偏移相当于按实验室同时性而非车系爱因斯坦同时性设置钟。钟的正常走时速率与一组钟的同步约定是不同问题。
\end{solution}

\begin{exercise}\label{ex:m1propersegments}
粒子先以$v_1$运动实验室时间$T_1$，再以$v_2$运动$T_2$，方向任意，变速时间忽略。求总固有时；若已知两端点坐标差，何时可以用端点间隔代替？
\Diagram{fig:m1segmentsq}{折线世界线的两段固有时须相加}{\STAxes\draw[very thick](0,0)--(2,2.5)--(1,4);\draw[guide](0,0)--(1,4);\node[right]at(2,2.5){$B$};\node[right]at(1,4){$C$};}
\end{exercise}
\begin{solution}
图\ref{fig:m1segmentsq}每段匀速，故$\tau=T_1\sqrt{1-v_1^2/c^2}+T_2\sqrt{1-v_2^2/c^2}$。端点间隔给的是直接连接两端的匀速世界线固有时$\tau_{AC}=\sqrt{(T_1+T_2)^2-|\vct v_1T_1+\vct v_2T_2|^2/c^2}$。
\Diagram{fig:m1segmentsa}{在端点共同静止系中，直线停留的钟积累时间最多}{\draw[axis](-1.5,0)--(2,0)node[right]{$x_*$};\draw[axis](0,0)--(0,4.3)node[above]{$ct_*$};\draw[very thick](0,0)--(1.3,2.1)--(0,3.8);\draw[guide](0,0)--(0,3.8);}
在图\ref{fig:m1segmentsa}的端点静止系，实际轨迹$\int\sqrt{1-v_*^2/c^2}\dd t_*\leq\int\dd t_*=\tau_{AC}$，等号要求全程$v_*=0$，即两段本来就是同一匀速速度矢量。仅两段速率相同而方向改变，仍不能用端点间隔代替曲线积分。
\end{solution}

\begin{exercise}\label{ex:m1inversevel}
实验室测得粒子$v_x=.8c,v_y=.3c$。求使其纵向速度为零的$x$向惯性系，以及该系中的横向速度和洛伦兹因子。
\Diagram{fig:m1inverseq}{选择变换速度消去纵向分量}{\draw[axis](0,0)--(4,0)node[right]{$x$};\draw[axis](0,0)--(0,2.3)node[above]{$y$};\draw[->,thick](0,0)--(3.2,1.2);\draw[guide](3.2,0)--(3.2,1.2);}
\end{exercise}
\begin{solution}
图\ref{fig:m1inverseq}中，设所求参考系沿$x$以$V$运动。速度变换的分母$1-Vv_x/c^2>0$，所以$v'_x=0$等价于$V=v_x=.8c$，其洛伦兹因子为$\gamma=5/3$。横向速度为
\[
 v'_y=\frac{v_y}{\gamma(1-Vv_x/c^2)}
 =\frac{.3c}{(5/3)(1-.64)}=.5c.
\]
因此新系中粒子仅沿$y'$运动，$\gamma_{v'}=(1-.25)^{-1/2}=2/\sqrt3$。
\Diagram{fig:m1inversea}{新系中轨迹纯横向，速率仍小于$c$}{\draw[axis](0,0)--(3.7,0)node[right]{$x'$};\draw[axis](0,0)--(0,2.7)node[above]{$y'$};\draw[very thick,->](0,0)--(0,2.1)node[right]{$.5c$};}
横向速度增大的原因可由微小事件差看出：$\dd y'=\dd y$，而
\[
 \dd t'=\gamma(1-Vv_x/c^2)\dd t=.6\dd t.
\]
相同横向位移除以更短的新系坐标时间，得到图\ref{fig:m1inversea}的$.5c$，并非横向位移本身被放大。四速度提供独立核对：原速率平方$.73c^2$，故
\[
 \gamma_v=\frac1{\sqrt{.27}},\qquad
 \gamma_{v'}=\gamma\gamma_v(1-.64)
 =\frac{.6}{\sqrt{.27}}=\frac2{\sqrt3}.
\]
只消去纵向速度仍未使粒子静止。若再在$S'$中沿$y'$以$.5c$变换，才到达粒子的静止系；两次非共线变换的空间轴关系须另行确定。
\end{solution}

\begin{exercise}\label{ex:m1numberangular}
各向同性光源在固有时内每秒发射$\dot N_0$个光子。求实验室按发射时间定义的角光子率，积分求总光子率，并与总功率不变比较。
\Diagram{fig:m1numberq}{同一批光子的数目不变，角密度与时间率改变}{\draw[axis](0,0)--(3.8,0)node[right]{$V$};\draw[curve](0,0)--(3,1);\draw[curve](0,0)--(3,1.5);\node[right]at(3,1.5){$\dd\Omega$};}
\end{exercise}
\begin{solution}
图\ref{fig:m1numberq}中，源系在$\dd\tau$内向$\dd\Omega'$发出的光子数为$\dd N=\dot N_0\dd\tau\dd\Omega'/(4\pi)$。令$\delta=[\gamma(1-\beta\cos\theta)]^{-1}$。由光行差，
\[
 \frac{\dd\cos\theta}{\dd\cos\theta'}
 =\frac1{\gamma^2(1+\beta\cos\theta')^2}=\delta^{-2},
\]
方位角不变，所以$\dd\Omega=\delta^{-2}\dd\Omega'$。再用实验室发射时间$\dd t=\gamma\dd\tau$，得到
\[
 \frac{\dd\dot N}{\dd\Omega}=\frac{\dot N_0}{4\pi\gamma}\delta^2.
\]
\Diagram{fig:m1numbera}{单次光子数不变，时间膨胀降低实验室总发射率}{\draw[axis](0,0)--(4.2,0)node[right]{$\beta$};\draw[axis](0,0)--(0,2.7)node[above]{$\dot N/\dot N_0$};\draw[curve,domain=0:1,samples=70]plot({3.5*\x},{2.3*sqrt(1-\x*\x)});\node[left]at(0,2.3){$1$};}
全角积分可以直接计算。以$\mu=\cos\theta$，
\begin{align*}
 \int\delta^2\dd\Omega
 &=\frac{2\pi}{\gamma^2}\int_{-1}^1\frac{\dd\mu}{(1-\beta\mu)^2}\\
 &=\frac{2\pi}{\gamma^2}
 \left[\frac{1}{\beta(1-\beta\mu)}\right]_{-1}^1\\
 &=\frac{2\pi}{\beta\gamma^2}
 \left[\frac1{1-\beta}-\frac1{1+\beta}\right]=4\pi.
\end{align*}
在$\beta=0$处取连续极限。因此图\ref{fig:m1numbera}给出总光子率$\dot N=\dot N_0/\gamma$。若源系每光子能量为$E_0$，同一批各向同性光子的总动量为零，整体能量变换便给平均能量$\langle E\rangle=\gamma E_0$。于是$P=\dot N\langle E\rangle=\dot N_0E_0$。按实验室发射时间计，光子数率因时间膨胀减小，平均单光子能量增加，两者相抵。接收时间率还须计入传播延迟。
\end{solution}

\begin{exercise}\label{ex:m1redblue}
单色运动源静止频率$\nu_0$。求实验室观察到无频移的方向，并确定蓝移、红移的角域。
\Diagram{fig:m1redq}{多普勒因子随实验室观察角变化}{\draw[axis](0,0)--(4,0)node[right]{$V$};\draw[curve](0,0)--(2.3,1.6);\draw[curve](0,0)--(-1,1.6);\draw[guide](0,0)--(0,2.2);\node at(1,.35){$\theta$};}
\end{exercise}
\begin{solution}
图\ref{fig:m1redq}中，$0<\beta<1$时$\nu/\nu_0=[\gamma(1-\beta\cos\theta)]^{-1}$。令此比为1，得到
\[
 \cos\theta_* =\frac{1-1/\gamma}{\beta}
 =\frac{\gamma-1}{\gamma\beta}.
\]
利用$(\gamma-1)(\gamma+1)=\gamma^2\beta^2$，可改写为$\cos\theta_*=\gamma\beta/(\gamma+1)$。该值介于0和1之间，对应唯一的前向无频移圆锥。
\Diagram{fig:m1reda}{无频移锥内蓝移，锥外红移}{\draw[axis](-1,0)--(4,0)node[right]{$x$};\draw[curve](0,0)--(1.8,2);\draw[curve](0,0)--(1.8,-2);\node[right]at(2,.5){蓝移};\node[left]at(-.2,1.5){红移};\node[above]at(1.8,2){$\theta_*$};}
图\ref{fig:m1reda}的锥内满足$\cos\theta>\cos\theta_*$，多普勒分母小于1，因而蓝移；锥外红移。将边界方向反变换到源系，有
\[
 \cos\theta'_* =\frac{\cos\theta_*-\beta}{1-\beta\cos\theta_*}
 =-\frac{\gamma\beta}{\gamma+1}=-\cos\theta_*.
\]
源系略向后的发射方向也可能在实验室无频移。低速展开$\gamma=1+\beta^2/2+\cdots$给$\cos\theta_*\simeq\beta/2$，边界略小于$\pi/2$；高速时$\cos\theta_*\to1$。当$\beta=0$时所有方向都无频移，不能把$\beta\to0$所得的边界锥误认为静止光源只有一个无频移方向。
\end{solution}

